\section{Il linguaggio Scala: programmazione funzionale applicata ai Big Data}\label{sec:scala}

Scala è un linguaggio multi-paradigma che combina programmazione orientata agli oggetti e programmazione funzionale sulla Java Virtual Machine (JVM).
La sua progettazione punta a un type system espressivo, a immutabilità per default nelle strutture dati e a un'interoperabilità diretta con l'ecosistema Java, caratteristiche che lo rendono adatto a sistemi distribuiti di medie e grandi dimensioni.

Nel contesto di questo lavoro, Scala non è stato scelto per mere preferenze sintattiche, bensì perché rappresenta il \emph{linguaggio nativo} di Apache Spark.
Le API di Spark sono scritte in Scala; le trasformazioni su \texttt{Dataset} e \texttt{DataFrame} sfruttano encoder derivati dal type system, evitando la serializzazione Python--JVM che caratterizza PySpark e che introduce overhead non trascurabile in pipeline ad alta frequenza~\cite{odersky2008programming,zaharia2016apache}.

Il progetto adotta la versione \textbf{2.13.18}, come definito in \texttt{build.sbt}, con prefisso di package \texttt{it.utiu.thesis}.
Tale scelta consente di sfruttare le ottimizzazioni della serie 2.13 (collezioni più efficienti, miglioramenti al compilatore) mantenendo compatibilità con Spark 4.1.1.

\subsection{Immutabilità e composizione}

Un principio ricorrente nel codice sorgente è l'uso di trasformazioni \emph{lazy} su DataFrame anziché mutazioni esplicite dello stato.
In \texttt{EcgModelTrainer}, la pipeline ML è costruita concatenando \texttt{StringIndexer}, \texttt{VectorAssembler} e \texttt{RandomForestClassifier} in un unico oggetto \texttt{Pipeline} immutabile.
Questo stile riduce la superficie di errori in ambienti concorrenti: il thread del predictor streaming e quello del simulatore operano su istanze Spark distinte senza condividere stato mutabile.

\subsection{Confronto con alternative}

PySpark offre una curva di apprendimento più rapida, ma ogni chiamata attraversa il confine Python--JVM tramite Py4J.
Per un simulatore che invia un messaggio al secondo il costo è trascurabile; per un flusso clinico reale con centinaia di eventi al minuto per paziente, la latenza di serializzazione diventa rilevante.
Java puro, d'altra parte, manca della concisione funzionale necessaria a esprimere trasformazioni su dataset in poche righe.
Scala occupa una posizione intermedia favorevole: performance JVM, sintassi adatta al data engineering, integrazione diretta con Spark.

\section{Apache Spark: Core, Spark SQL e MLlib}\label{sec:spark}

Apache Spark è un motore di calcolo distribuito che generalizza il modello MapReduce introducendo operazioni in-memory su strutture resilienti~\cite{zaharia2012resilient,zaharia2016apache}.
Il progetto utilizza la versione \textbf{4.1.1}, con moduli \texttt{spark-core}, \texttt{spark-sql}, \texttt{spark-mllib} e connettore \texttt{spark-sql-kafka-0-10}.

\subsection{Modello di esecuzione}

Spark organizza il lavoro attorno a un \emph{driver} (processo coordinatore) e a uno o più \emph{executor} (processi worker).
Il driver traduce le operazioni su DataFrame in un DAG (Directed Acyclic Graph) di stage, ciascuno composto da task parallelizzabili.
Nel setup sperimentale di questo lavoro, il master è configurato come \texttt{local[*]} in \texttt{Config.scala}: tutti i core della macchina host vengono impiegati su un singolo nodo.
Questa configurazione è adeguata alla prototipazione accademica, ma non rappresenta un deployment produttivo su cluster.

\subsection{Spark SQL e Catalyst}

Spark SQL fornisce un'API dichiarativa su tabelle strutturate.
L'ottimizzatore Catalyst applica regole di rewrite, predicate pushdown e codegen Tungsten, generando bytecode JVM specializzato per le operazioni ripetitive (filtri, proiezioni, join).
Nel trainer, la lettura CSV con \texttt{inferSchema=true} e la successiva \texttt{groupBy("type").count()} beneficiano di questo pipeline di ottimizzazione senza intervento manuale dell'implementatore.

\subsection{MLlib e la Pipeline API}

MLlib espone algoritmi di machine learning distribuiti attraverso un'astrazione a stadi: \texttt{Estimator} (addestrabili), \texttt{Transformer} (applicabili) e \texttt{Pipeline} (composizione sequenziale)~\cite{meng2016mllib}.
Il vantaggio principale è la \textbf{serializzazione atomica}: salvando \texttt{pipelineModel.write().save(...)} si persistono contemporaneamente l'indicizzatore delle etichette, l'assemblatore delle feature e il classificatore Random Forest.
In fase di serving, \texttt{PipelineModel.load(...)} ricostruisce l'intera catena di trasformazione, garantendo coerenza tra training e inferenza.

\subsection{Structured Streaming}

Structured Streaming estende il modello batch di Spark SQL a flussi illimitati, trattando ogni micro-batch come una ``tabella unbounded'' aggiornata nel tempo~\cite{armbrust2018structured}.
L'API è dichiarativa: \texttt{readStream} definisce la sorgente, le trasformazioni su DataFrame replicano la logica batch, \texttt{writeStream} configura il sink.
Il motore gestisce offset, checkpoint e recovery in modo trasparente, riducendo la complessità rispetto alle API Spark Streaming basate su DStream (discontinuate).

Nel componente \texttt{EcgStreamingPredictor}, Structured Streaming consuma messaggi Kafka, applica il modello ML e scrive su Delta Lake tramite \texttt{foreachBatch}.
La scelta di \texttt{foreachBatch} rispetto a un sink nativo consente logica custom (persist, controllo \texttt{isEmpty}, logging) a costo di una gestione più esplicita del ciclo di vita del batch.

\section{Apache Kafka: architettura publish-subscribe e disaccoppiamento}\label{sec:kafka}

Apache Kafka è un sistema di messaggistica distribuito organizzato attorno a \emph{topic}, \emph{partizioni} e \emph{offset}~\cite{kreps2011kafka}.
I producer scrivono record append-only; i consumer leggono attraverso consumer group con bilanciamento automatico del carico.

\subsection{Architettura log-centrica}

Kafka non è una semplice coda di messaggi: ogni topic è un log commitato replicato, consultabile con semantica di offset.
Questo modello si adatta bene a pipeline di analytics dove più consumer possono rileggere lo stesso flusso a diversi ritmi (real-time vs batch di riprocessamento).

Nel progetto, il topic \texttt{ecg-input-stream} funge da buffer tra il simulatore (\texttt{EcgPatientSimulator}) e il predictor (\texttt{EcgStreamingPredictor}).
Il disaccoppiamento temporale è evidente: il producer può inviare un battito al secondo mentre Spark elabora micro-batch con latenza variabile, senza che le due velocità debbano coincidere.

\subsection{Modalità KRaft e deprecazione di ZooKeeper}

Storicamente, Kafka dipendeva da Apache ZooKeeper per la gestione dei metadati del cluster (broker, topic, partition leader).
A partire dalla versione 3.x, il progetto Apache ha introdotto KRaft (Kafka Raft): un quorum interno che elimina la dipendenza esterna~\cite{apache2023kraft}.

Il file \texttt{docker-compose.yml} del repository configura un broker in modalità combinata:

\begin{itemize}
  \item \texttt{KAFKA\_PROCESS\_ROLES=broker,controller} --- un singolo nodo svolge entrambi i ruoli;
  \item \texttt{KAFKA\_CONTROLLER\_QUORUM\_VOTERS=1@localhost:9093} --- quorum a nodo singolo per sviluppo locale;
  \item \texttt{KAFKA\_OFFSETS\_TOPIC\_REPLICATION\_FACTOR=1} --- coerente con un cluster di prova senza ridondanza.
\end{itemize}

In produzione, i ruoli broker e controller andrebbero separati su macchine distinte con replication factor $\geq 3$.
Per la tesi, la configurazione minimalista è sufficiente a dimostrare l'integrazione senza introdure complessità operativa non necessaria.

\subsection{Garanzie di consegna}

Kafka offre tre livelli di garanzia: at-most-once, at-least-once ed exactly-once (con transazioni e idempotenza del producer).
Spark Structured Streaming, combinato con checkpoint e sink idempotente (Delta Lake), mira a exactly-once \emph{a livello di micro-batch}: ogni batch viene processato una sola volta anche in caso di restart del driver.

La Tabella~\ref{tab:kraft-vs-zk} riassume il confronto tra l'architettura classica e KRaft.

\begin{table}[H]
  \centering
  \caption{Confronto tra Kafka con ZooKeeper e Kafka in modalità KRaft.}
  \label{tab:kraft-vs-zk}
  \begin{tabular}{@{}lll@{}}
    \toprule
    \textbf{Aspetto} & \textbf{ZooKeeper} & \textbf{KRaft} \\
    \midrule
    Componenti esterni & Sì (cluster ZK separato) & No \\
    Tempo di failover & Secondi--decine di s & Sub-secondo (target) \\
    Complessità operativa & Alta & Ridotta \\
    Scalabilità metadata & Limitata & Migliorata \\
    \bottomrule
  \end{tabular}
\end{table}

\section{Delta Lake: il paradigma Data Lakehouse}\label{sec:delta-lake}

Delta Lake estende lo storage columnar (Parquet) con un \emph{transaction log} che garantisce proprietà ACID su object store e filesystem distribuiti~\cite{armbrust2020delta}.
Il concetto di \emph{Lakehouse} unifica la flessibilità del data lake (formati aperti, costi contenuti) con le garanzie transazionali tipiche dei data warehouse.

\subsection{Limiti di Parquet puro}

Parquet da solo non supporta transazioni concorrenti: due writer che appendono contemporaneamente possono corrompere la directory.
Non esiste versioning nativo, né meccanismo di rollback.
Per una pipeline di predizioni cliniche, dove ogni battito classificato deve essere persistito in modo affidabile, queste lacune sono inaccettabili.

\subsection{Transaction log e ACID}

Delta Lake mantiene una directory \texttt{\_delta\_log} con file JSON (e checkpoint Parquet) che registrano ogni operazione: append, overwrite, merge, delete.
Le scritture sono atomiche: o un batch è visibile per intero, o non lo è affatto.
Il meccanismo di \emph{optimistic concurrency control} risolve conflitti tra writer concorrenti.

\subsection{Integrazione con Spark}

Il progetto configura Spark con:

\begin{lstlisting}[language=Scala,caption={Configurazione Delta Lake in SparkSession.},label={lst:delta-config}]
.config("spark.sql.extensions",
  "io.delta.sql.DeltaSparkSessionExtension")
.config("spark.sql.catalog.spark_catalog",
  "org.apache.spark.sql.delta.catalog.DeltaCatalog")
\end{lstlisting}

La stessa configurazione compare in \texttt{EcgModelTrainer} e \texttt{EcgStreamingPredictor}, anche se solo il predictor scrive effettivamente su Delta.
La uniformità evita sorprese se in futuro il trainer dovesse salvare metriche di validazione in formato Delta.

I risultati delle predizioni confluiscono in \texttt{spark-warehouse/ecg\_predictions}, con checkpoint separato in \texttt{spark-warehouse/checkpoints/ecg\_predictions}.
Questa separazione è una buona pratica: il checkpoint contiene metadati di offset Spark, non dati di business.

\section{Sintesi: requisiti e tecnologie adottate}\label{sec:sintesi-stack}

La Tabella~\ref{tab:requisiti-tecnologie} collega i requisiti emersi nell'introduzione alle scelte tecnologiche effettuate.

\begin{longtable}{@{}p{4.5cm}p{4cm}p{5.5cm}@{}}
  \caption{Mappatura requisiti--tecnologie.} \label{tab:requisiti-tecnologie} \\
  \toprule
  \textbf{Requisito} & \textbf{Tecnologia} & \textbf{Meccanismo} \\
  \midrule
  \endfirsthead
  \multicolumn{3}{c}{\tablename\ \thetable\ --- \emph{continua}} \\
  \toprule
  \textbf{Requisito} & \textbf{Tecnologia} & \textbf{Meccanismo} \\
  \midrule
  \endhead
  Disaccoppiamento ingestion/processing & Kafka & Topic \texttt{ecg-input-stream} con offset indipendenti \\
  Training su dati storici & Spark MLlib & \texttt{EcgModelTrainer}, split 80/20, \texttt{Pipeline.fit} \\
  Inferenza near-real-time & Structured Streaming & Micro-batch da Kafka, \texttt{pipelineModel.transform} \\
  Persistenza affidabile & Delta Lake & Append transazionale, \texttt{foreachBatch} \\
  Fault tolerance & Checkpoint Spark & \texttt{STREAMING\_CHECKPOINT\_PATH} \\
  Linguaggio unificato & Scala 2.13 & Batch e streaming nello stesso codebase \\
  \bottomrule
\end{longtable}

Il capitolo successivo traduce queste scelte in un'architettura concreta, descrivendo il flusso dei dati e il ruolo di ciascun modulo (\texttt{EcgModelTrainer}, \texttt{EcgPatientSimulator}, \texttt{EcgStreamingPredictor}) orchestrato da \texttt{AppRunner}.
